Pudo verificarse el cumplimiento de la ley de Vegard para soluciones sólidas
$\ce{KCl_{1-x}Br_x}$ al observarse una relación lineal entre los parámetros de red
de dichas soluciones.

Además pudo determinarse los parámetros de red de diferentes compuestos con errores menores
al $0,08\%.$ 

Por otro lado, mediante el uso de difractogramas, se pudo analizar la estructura de diferentes compósitos,
como es el caso de la sal light y del compósito de $\ce{KCl}$ y $\ce{KBr}$ con fracciones molares de $0,4$ y $0,6$ respectivamente.

Luego mediante el uso de un microscopio electrónico de barrido se observó la morfología de las soluciones de $\ce{KCl_{1-x}Br_x},$
observando una mejor cristalización en la muestra que se dejó más tiempo cristalizando.

Finalmente se comprobó que las estructuras de tipo sal de roca son estructuras compactas.